\cleardoublepage
\chapter{PENDAHULUAN}
\label{chap:pendahuluan}

\section{Latar Belakang}

Kegiatan campaign pada industri perbankan membutuhkan data yang dapat digunakan untuk menentukan target, menjalankan komunikasi, dan membaca hasil setelah campaign berjalan. Pekerjaan tersebut tidak berhenti pada pembuatan query. Data perlu dipahami berdasarkan konteks bisnis, diperiksa kualitasnya, diproses melalui alur yang dapat dijalankan ulang, lalu disajikan dalam bentuk yang dapat dibaca oleh unit terkait.

Selama Kerja Praktik di PT Bank Syariah Indonesia (Persero) Tbk, penulis terlibat pada pekerjaan yang berada di antara kebutuhan bisnis dan implementasi data. Penulis mempelajari alur data untuk campaign, mengembangkan query SQL, menguji proses ETL menggunakan Talend, menyusun analisis deskriptif campaign financing, serta mengembangkan spatial clustering dari data log aplikasi digital.

Tiga jenis pekerjaan tersebut dipilih sebagai isi utama laporan karena mewakili alur kerja yang penulis jalankan selama Kerja Praktik. Pekerjaan dimulai dari memahami kebutuhan, mengambil dan menyiapkan data, memeriksa hasil, lalu menyusun analisis yang dapat digunakan sebagai bahan diskusi bisnis. Fokus laporan berada pada proses dan keputusan teknis. Detail internal yang tidak diperlukan untuk menjelaskan metode tidak dicantumkan.

\section{Tempat dan Waktu Pelaksanaan}

Kerja Praktik dilaksanakan di PT Bank Syariah Indonesia (Persero) Tbk pada fungsi Campaign \& Leads Management. Lokasi kerja yang tercantum pada dokumen pelaksanaan Kerja Praktik berada di Jl. Cisitu Baru, Dago, Coblong, Bandung, Jawa Barat 40135.

Pelaksanaan Kerja Praktik berlangsung pada periode Juli sampai September 2026. Tanggal mulai dan selesai secara administratif akan disesuaikan kembali dengan dokumen final perusahaan sebelum laporan dikumpulkan.

\section{Ruang Lingkup Kerja Praktik}

Ruang lingkup pekerjaan dalam laporan ini dibagi menjadi tiga bagian.

\begin{enumerate}
    \item Pengembangan query SQL dan proses ETL untuk kebutuhan data campaign.
    \item Analisis deskriptif efektivitas campaign financing pada tingkat agregat.
    \item Spatial clustering dan visualisasi data log aplikasi digital untuk membaca pola lokasi.
\end{enumerate}

Laporan ini tidak mencantumkan data mentah, identifier nasabah, nomor telepon, nama tabel internal, nama kolom sensitif, query produksi secara lengkap, kode campaign, maupun nilai bisnis internal yang bersifat terbatas. Nama produk atau unit hanya digunakan jika diperlukan untuk menjelaskan konteks secara umum. Angka operasional yang tidak diperlukan untuk menjelaskan metode juga tidak ditampilkan.

\section{Sistematika Penulisan}

Bab I menjelaskan latar belakang, tempat dan waktu pelaksanaan, ruang lingkup, serta batasan penulisan. Bab II menjelaskan profil perusahaan dan unit kerja pada tingkat yang relevan terhadap kegiatan Kerja Praktik. Bab III menjelaskan pekerjaan yang dilakukan, mulai dari pengolahan data, ETL, analisis campaign, sampai spatial clustering. Bab IV berisi rangkuman pekerjaan, hasil, evaluasi diri, dan saran.
